\documentclass{article}
\usepackage[utf8]{inputenc}
\usepackage{amsmath}
\usepackage{geometry}
\geometry{margin=2cm}
\begin{document}

\section*{Questão 117 — ENEM 2024 — Ciências da Natureza}

\subsection*{Interpretação e estratégia}

A questão apresenta a situação de uma ambulância em alta velocidade se aproximando de um radar que emite ondas de rádio de frequência \(f_0\). Essas ondas são refletidas pela ambulância e recebidas com frequência \(f_r\). O operador do radar percebe o som da sirene da ambulância de forma alterada devido ao movimento.

\textbf{Termos da questão:}
\begin{itemize}
  \item \textit{Percepção}: Como o operador percebe a mudança do som da sirene na aproximação.
  \item \textit{Relação entre frequências}: Comparação entre \(f_r\) e \(f_0\).
\end{itemize}

\textbf{Objetivo:} Descrever a percepção sonora do operador em relação à sirene e determinar a relação entre as frequências \(f_r\) e \(f_0\) durante a aproximação.

\textbf{Estratégia adotada:} Aplicar o efeito Doppler para ondas refletidas e relacionar com a percepção auditiva do operador.

\subsection*{Conteúdo e mini revisão}

\begin{tabular}{|p{0.3\textwidth}|p{0.4\textwidth}|p{0.3\textwidth}|}
\hline
\textbf{Conceito} & \textbf{Ideia chave} & \textbf{Aplicação no problema} \\
\hline
Efeito Doppler & Movimento relativo altera frequência percebida & Durante a aproximação, \(f_r > f_0\) \\
\hline
Frequência e som & Frequência maior significa som mais agudo & O operador ouve som mais agudo da sirene na aproximação \\
\hline
Ondas refletidas pelo radar & Ondas de rádio refletidas sofrem variação de frequência & \(f_r\) é alterada pela velocidade da ambulância, \(f_r > f_0\) na aproximação \\
\hline
\end{tabular}

\textbf{Teoria:}

O efeito Doppler ocorre quando há movimento relativo entre fonte e observador, alterando a frequência percebida. Para ondas sonoras, aproximação eleva a frequência e, para ondas refletidas de radar, a alteração também acontece, revelando a velocidade da ambulância. 

\subsection*{Resolução e pulo do gato}

A ambulância emite som pela sirene (freqüência \(f_0\)) e está se deslocando em direção ao radar. O radar emite ondas de rádio com frequência \(f_0\) que são refletidas pela ambulância. 

Pelo efeito Doppler, a frequência da onda refletida recebida pelo radar, \(f_r\), será maior do que \(f_0\) devido à aproximação da ambulância:

\[
f_r > f_0
\]

Simultaneamente, o operador que escuta o som da sirene percebe este som mais agudo, pois a frequência do som também aumenta na aproximação.

\textbf{Pulo do gato:} Na aproximação, a frequência das ondas refletidas aumenta (\(f_r > f_0\)) e o som fica mais agudo, confirmando o efeito Doppler clássico.

\subsection*{Armadilhas e código da questão}

\textbf{Armadilhas:}
\begin{itemize}
  \item Achar que frequência diminui na aproximação (som mais grave).
  \item Confundir frequência das ondas refletidas com a frequência sonora original.
  \item Ignorar o efeito Doppler e pensar em frequência constante.
  \item Interpretar som mais agudo como frequência menor.
\end{itemize}

\textbf{Código da questão:} CN\_EFEITODOPPLER\_BASICO\_001

\bigskip

\noindent{\small\emph{Esta resolução foi elaborada com base na análise técnica e didática da questão, garantindo clareza e coerência conceitual no ensino do efeito Doppler.}}

\end{document}